\documentclass[12pt,a4paper]{article}

% === МОНГОЛ КИРИЛЛ ҮСГИЙН ДЭМЖЛЭГ (pdfLaTeX-д зөв арга) ===
\usepackage[T2A,T1]{fontenc}        % T2A = Кирилл үсэг
\usepackage[mongolian,english]{babel}  % mongolian = Монгол кирилл

% === Бусад хэрэгтэй багцууд ===
\usepackage{booktabs}
\usepackage{caption}
\usepackage{geometry}
\geometry{margin=2.5cm}

\title{Түүврийн баримт бичиг}
\author{}
\date{\today}

\begin{document}

\maketitle

\section{Оршил}

Монгол Улс нь дотоодын нийт бүтээгдэхүүн (ДНБ)-ээрээ дэлхийд дунд зэрэгт ордог. Газрын гадаргаа нь 1.5 сая км² талбайтай, хүн ам нь 3.5 сая орчим.

\section{Өгөгдөл}

\begin{table}[h]
\centering
\begin{tabular}{lcr}
\toprule
\textbf{Бүтээгдэхүүн} & \textbf{Тоо хэмжээ} & \textbf{Үнэ} \\
\midrule
А ангилал & 10 & 100 \\
Б ангилал & 20 & 200 \\
В ангилал & 30 & 300 \\
Г ангилал & 40 & 400 \\
\bottomrule
\end{tabular}
\caption{Жишээ хүснэгт}
\end{table}

\section{Дүгнэлт}

Монгол Улс цаашид уул уурхай, хөдөө аж ахуй, аялал жуулчлалын салбараа хөгжүүлэх боломжтой.

\end{document}